% !TEX program = xelatex
% Lernplatt - by Can Siebert
% Kurs 71 (Dig 42) - Tag 15 - Aufgabenheft
% Change Leadership: Rollenklaerung, Stakeholder/Kommunikation, Widerstand produktiv, Befaehigen & Verankern
%
% Self-contained xelatex source. Kein biblatex, keine externen Bilder.

\documentclass[11pt,a4paper,oneside]{article}

\usepackage{polyglossia}
\setdefaultlanguage{german}

\usepackage[a4paper,margin=2.5cm]{geometry}

\usepackage{fontspec}
\defaultfontfeatures{Ligatures=TeX}

\IfFontExistsTF{Charter}{%
  \setmainfont{Charter}%
}{%
  \IfFontExistsTF{Noto Serif}{%
    \setmainfont{Noto Serif}%
  }{%
    \setmainfont{Liberation Serif}%
  }%
}

\IfFontExistsTF{Source Sans 3}{%
  \setsansfont{Source Sans 3}%
}{%
  \IfFontExistsTF{Source Sans Pro}{%
    \setsansfont{Source Sans Pro}%
  }{%
    \setsansfont{Liberation Sans}%
  }%
}

\newcommand{\caveatfontfallback}{\itshape}
\IfFontExistsTF{Caveat}{%
  \newfontfamily\caveatfont{Caveat}[Scale=1.15]%
}{%
  \newcommand\caveatfont{\caveatfontfallback}%
}

\setmonofont{Liberation Mono}

\usepackage{microtype}
\linespread{1.15}

\usepackage{xcolor}
\definecolor{lpInk}{RGB}{20,28,38}
\definecolor{lpAccent}{RGB}{157,74,26}
\definecolor{lpAccentLight}{RGB}{246,228,212}
\definecolor{lpAmber}{RGB}{222,138,30}
\definecolor{lpAmberLight}{RGB}{255,243,217}
\definecolor{lpAmberBorder}{RGB}{196,118,18}
\definecolor{lpGray}{RGB}{96,104,114}
\definecolor{lpGrayLight}{RGB}{238,240,243}
\definecolor{lpGreen}{RGB}{34,120,72}
\definecolor{lpGreenLight}{RGB}{222,238,228}
\definecolor{lpGreenBorder}{RGB}{24,96,58}

\definecolor{lpL1}{RGB}{34,120,72}
\definecolor{lpL2}{RGB}{22,90,140}
\definecolor{lpL3}{RGB}{170,42,52}

\usepackage[most]{tcolorbox}
\tcbuselibrary{breakable,skins}

\newtcolorbox{infobox}[1][Hinweis]{%
  enhanced, breakable,
  colback=lpAccentLight,
  colframe=lpAccent,
  boxrule=0.6pt,
  arc=2pt,
  left=10pt,right=10pt,top=8pt,bottom=8pt,
  fonttitle=\sffamily\bfseries\color{white},
  coltitle=white,
  title={#1},
  attach boxed title to top left={xshift=8pt,yshift=-8pt},
  boxed title style={size=small, colback=lpAccent, colframe=lpAccent, boxrule=0.4pt, arc=1pt},
  top=14pt
}

\newtcolorbox{antwortbox}[1][Ihre Antwort]{%
  enhanced, breakable,
  colback=white,
  colframe=lpGray,
  boxrule=0.4pt,
  arc=2pt,
  left=10pt,right=10pt,top=8pt,bottom=8pt,
  fonttitle=\sffamily\bfseries\color{white},
  coltitle=white,
  title={#1},
  attach boxed title to top left={xshift=8pt,yshift=-8pt},
  boxed title style={size=small, colback=lpGray, colframe=lpGray, boxrule=0.4pt, arc=1pt},
  top=14pt
}

\usepackage{enumitem}
\setlist{itemsep=2pt, topsep=4pt, parsep=0pt}
\setlist[itemize,1]{label={\textbullet},leftmargin=1.6em}
\setlist[itemize,2]{label={\textendash},leftmargin=1.4em}
\setlist[enumerate,1]{label=\arabic*., leftmargin=1.8em}

\usepackage{tabularx}
\usepackage{array}
\usepackage{booktabs}
\usepackage{longtable}

\usepackage{fancyhdr}
\usepackage{lastpage}
\pagestyle{fancy}
\fancyhf{}
\renewcommand{\headrulewidth}{0.3pt}
\renewcommand{\footrulewidth}{0pt}
\fancyhead[L]{\footnotesize\sffamily\color{lpGray}Aufgabenheft \textperiodcentered{} Kurs 71 \textperiodcentered{} Tag 15}
\fancyhead[R]{\footnotesize\sffamily\color{lpGray}Change Leadership \& Verankerung}
\fancyfoot[L]{\footnotesize\sffamily\color{lpGray}\textcopyright{} Can Siebert}
\fancyfoot[R]{\footnotesize\sffamily\color{lpGray}\thepage{} / \pageref{LastPage}}

\fancypagestyle{plain}{%
  \fancyhf{}%
  \renewcommand{\headrulewidth}{0pt}%
  \fancyfoot[C]{\footnotesize\sffamily\color{lpGray}\thepage{} / \pageref{LastPage}}%
}

\usepackage[explicit]{titlesec}

\titleformat{\section}[block]
  {\sffamily\Large\bfseries\color{lpAccent}}
  {\thesection.}{0.6em}{#1}
  [{\vspace{2pt}\color{lpAccent}\titlerule[0.6pt]}]

\titleformat{\subsection}[block]
  {\sffamily\large\bfseries\color{lpInk}}
  {\thesubsection}{0.6em}{#1}

\titlespacing*{\section}{0pt}{14pt}{8pt}
\titlespacing*{\subsection}{0pt}{10pt}{4pt}

\usepackage[hidelinks,unicode]{hyperref}

\newcommand{\akzent}[1]{{\caveatfont\color{lpAmber}#1}}
\newcommand{\fachbegriff}[1]{\textbf{#1}}

\newcommand{\niveaubadge}[2]{%
  \tcbox[on line, colback=#1, colframe=#1, coltext=white,
         boxrule=0pt, arc=2pt, left=5pt,right=5pt,top=1pt,bottom=1pt,
         nobeforeafter]{\sffamily\bfseries\footnotesize #2}%
}
\newcommand{\nivLi}{\niveaubadge{lpL1}{L1\,\textbullet\,Wissen}}
\newcommand{\nivLii}{\niveaubadge{lpL2}{L2\,\textbullet\,Anwendung}}
\newcommand{\nivLiii}{\niveaubadge{lpL3}{L3\,\textbullet\,Management}}

\newcommand{\zeit}[1]{%
  \tcbox[on line, colback=lpGrayLight, colframe=lpGrayLight, coltext=lpInk,
         boxrule=0pt, arc=2pt, left=5pt,right=5pt,top=1pt,bottom=1pt,
         nobeforeafter]{\sffamily\footnotesize\textbf{$\sim$\,#1}}%
}

\newcommand{\aufgabenkopf}[4]{%
  \par\vspace{6pt}
  \noindent{\sffamily\Large\bfseries\color{lpAccent}Aufgabe #1 \textendash{} #2}\par
  \vspace{2pt}
  \noindent{\color{lpAccent}\rule{\linewidth}{0.6pt}}\par
  \vspace{4pt}
  \noindent #3 \hfill \zeit{#4}\par
  \vspace{6pt}
}

\newcommand{\ytquery}[1]{%
  \par\vspace{4pt}
  \noindent{\footnotesize\sffamily\color{lpGray}\textbf{YouTube-Suchquery:} \texttt{#1}}\par
}

\newcommand{\antwortlinien}[1]{%
  \par\vspace{6pt}
  \foreach \x in {1,...,#1}{%
    \noindent\makebox[\linewidth]{\dotfill}\\[6pt]
  }
}
\usepackage{pgffor}

% =========================================================================
\begin{document}

\begin{titlepage}
  \pagestyle{empty}
  \vspace*{1.5cm}
  \noindent{\sffamily\footnotesize\color{lpGray}LERNPLATT \textperiodcentered{} BY CAN SIEBERT}\\[2pt]
  \noindent{\color{lpAccent}\rule{\linewidth}{1.2pt}}\\[4pt]
  \noindent{\sffamily\footnotesize\color{lpGray}AUFGABENHEFT \textperiodcentered{} MODUL 7 \textperiodcentered{} TAG 15 VON 16 \textperiodcentered{} KURS 71 (Dig 42) \textperiodcentered{} 18.\,Juni 2026}

  \vspace{3cm}
  {\sffamily\fontsize{30}{34}\selectfont\bfseries\color{lpInk}Aufgabenheft\\Tag 15\par}

  \vspace{0.7cm}
  {\sffamily\large\color{lpGray}Nachhaltige Prozess\-ver\"anderung\\f\"uhren \textendash{} Rollen, Widerstand, Befaehigung\\und Verankerung in elf Aufgaben.\par}

  \vspace{1.5cm}
  {\caveatfont\LARGE\color{lpAmber}11 Aufgaben \textperiodcentered{} L1 \textperiodcentered{} L2 \textperiodcentered{} L3}\par

  \vfill

  \noindent\begin{minipage}{0.55\linewidth}
    {\sffamily\footnotesize\color{lpGray}Niveau-Mix:}\\
    {\sffamily\small\color{lpInk}3\,\textsf{L1} \textperiodcentered{} 5\,\textsf{L2} \textperiodcentered{} 3\,\textsf{L3}}\\[10pt]
    {\sffamily\footnotesize\color{lpGray}Bearbeitung:}\\
    {\sffamily\small\color{lpInk}Selbstbearbeitung im Lernpfad}
  \end{minipage}\hfill
  \begin{minipage}{0.4\linewidth}
    \raggedleft
    {\sffamily\footnotesize\color{lpGray}Dozent:}\\
    {\sffamily\small\color{lpInk}Can Siebert}\\[10pt]
    {\sffamily\footnotesize\color{lpGray}Begleitend:}\\
    {\sffamily\small\color{lpInk}Lehrskript \& L\"osungsheft}
  \end{minipage}

  \vspace{0.5cm}
  \noindent{\color{lpAccent}\rule{\linewidth}{0.6pt}}
\end{titlepage}

\section*{So arbeiten Sie mit diesem Heft}
\thispagestyle{plain}

\noindent Dieses Aufgabenheft enth\"alt elf Aufgaben zu Tag 15 \textendash{} \emph{Nachhaltige Prozessver\"anderung f\"uhren}: Change-Rollen (Sponsor, Change Lead, Process Owner, Line Manager, Champions), Stakeholder- und Kommunikationsstrategie, Widerstand als Information lesen, Befaehigung und Sustainment-Mechanik. Bezugsmodelle: \emph{Kotter} (8-Schritte), \emph{Lewin} (Unfreeze-Change-Refreeze), \emph{ADKAR} (Awareness-Desire-Knowledge-Ability-Reinforcement).

\begin{itemize}
  \item \nivLi{} \textendash{} \fachbegriff{Wissen.} Rollen, Modelle, Widerstandstypen. 5\,--\,10 Minuten.
  \item \nivLii{} \textendash{} \fachbegriff{Anwendung.} Stakeholder-Mapping, Kommunikations\-paket, Sustainment-Design. 15\,--\,30 Minuten.
  \item \nivLiii{} \textendash{} \fachbegriff{Management.} Change-Portfolio, Transparenzregeln, Stoppkriterien. 30\,--\,45 Minuten.
\end{itemize}

\noindent Jede Aufgabe enth\"alt eine Niveau-Markierung, einen Zeit-Richtwert, den Aufgabentext, einen Antwort-Freiraum und eine \emph{YouTube-Suchquery}.

\begin{infobox}[Hinweis]
Die Musterl\"osungen finden Sie im separaten \emph{L\"osungsheft Tag 15}. \emph{Widerstand ist Information}: Interessen, Risiken, Identit\"at, \"Uberlastung, fehlendes Vertrauen. Wer Widerstand nur ``brechen'' will, l\"ost ihn nicht \textendash{} sondern produziert Schattenprozesse.
\end{infobox}

\clearpage

% =========================================================================
% L1 AUFGABEN
% =========================================================================
\section*{Niveau L1 \textendash{} Wissen \& Reproduktion}
\addcontentsline{toc}{section}{L1 -- Wissen}

% --- AUFGABE 1 -----------------------------------------------------------
% LERNPLATT_HILFSVIDEOS
\section*{Hilfsvideos zu diesem Tag}
\addcontentsline{toc}{section}{Hilfsvideos zu diesem Tag}
\thispagestyle{plain}

\noindent Die folgenden YouTube-Erkl\"arvideos vermitteln die Inhalte, die Sie zur Bearbeitung der Aufgaben ben\"otigen. Klicken Sie auf den Titel, um das Video direkt zu \"offnen; die vollst\"andige URL steht in Klammern.

\begin{description}[style=nextline,leftmargin=18pt,labelindent=0pt,itemsep=8pt]
  \item[1.~Kotter — Das 8-Stufen-Modell der Veränderung] \href{https://youtu.be/7yiI_n79Aj8}{\textcolor{lpAccent}{\nolinkurl{https://youtu.be/7yiI_n79Aj8}}}\\[2pt]
  {\footnotesize\color{lpGray}(URL: \nolinkurl{https://youtu.be/7yiI_n79Aj8})}
  \item[2.~Lewin — Unfreeze, Change, Refreeze] \href{https://youtu.be/zz7xzY8c6Vo}{\textcolor{lpAccent}{https://youtu.be/zz7xzY8c6Vo}}\\[2pt]
  {\footnotesize\color{lpGray}(URL: \nolinkurl{https://youtu.be/zz7xzY8c6Vo})}
  \item[3.~Widerstand gegen Veränderung produktiv nutzen] \href{https://youtu.be/9zTC87rzeI8}{\textcolor{lpAccent}{https://youtu.be/9zTC87rzeI8}}\\[2pt]
  {\footnotesize\color{lpGray}(URL: \nolinkurl{https://youtu.be/9zTC87rzeI8})}
  \item[4.~Stakeholder-Management — Einfluss, Betroffenheit, Einbindung] \href{https://youtu.be/M7tmdJBQcZY}{\textcolor{lpAccent}{https://youtu.be/M7tmdJBQcZY}}\\[2pt]
  {\footnotesize\color{lpGray}(URL: \nolinkurl{https://youtu.be/M7tmdJBQcZY})}
\end{description}
\vspace{0.6em}
\noindent{\footnotesize\color{lpGray}\textit{Tipp:} \"Offnen Sie das Video, lassen Sie es im Hintergrund laufen und arbeiten Sie parallel an der Aufgabe.}

\clearpage

\aufgabenkopf{1}{F\"unf Change-Rollen + Verantwortlichkeiten}{\nivLi}{8\,min}

\noindent Eine gelungene Prozessver\"anderung braucht mindestens f\"unf Rollen: \emph{Sponsor}, \emph{Change Lead}, \emph{Process Owner}, \emph{Line Manager}, \emph{Champions / Multiplikatoren}. Erfolgsfaktoren: klares ``Warum'', Beteiligung, Quick Wins, psychologische Sicherheit, sichtbare F\"uhrung, messbare Wirkung.

\medskip
\noindent\textbf{Arbeitsauftrag:}

\begin{enumerate}
  \item Beschreiben Sie f\"ur jede der f\"unf Rollen \textbf{drei Verantwortlichkeiten} (z.\,B.\ Sponsor: gibt Richtung, sichert Ressourcen, entscheidet bei Konflikten).
  \item Erkl\"aren Sie in einem Satz, was passiert, wenn die \emph{Sponsor-Rolle} unterbesetzt ist (typisch: ``Wir machen einen Change ohne echtes Mandat'').
  \item Nennen Sie \textbf{drei Anzeichen}, dass die \emph{Champions} nicht wirken (z.\,B.\ keine Feedbacksignale aus dem Team, Champions wechseln aus, Kollegen reden nicht mit ihnen \"uber den Change).
\end{enumerate}

\ytquery{change leadership roles sponsor change lead process owner champions Kotter}

\antwortlinien{14}

\clearpage

% --- AUFGABE 2 -----------------------------------------------------------
\aufgabenkopf{2}{Change-Modelle: Kotter, Lewin, ADKAR}{\nivLi}{10\,min}

\noindent Drei klassische Modelle leiten Change-Vorhaben:

\begin{itemize}
  \item \emph{Lewin}: \emph{Unfreeze} \textrightarrow{} \emph{Change} \textrightarrow{} \emph{Refreeze}. Konzentriert sich auf den Aggregatzustand der Organisation.
  \item \emph{Kotter}: \emph{8 Schritte} \textendash{} Dringlichkeit, F\"uhrungskoalition, Vision, Kommunikation, Hindernisse beseitigen, Quick Wins, Konsolidieren, Verankern.
  \item \emph{ADKAR}: \emph{Awareness, Desire, Knowledge, Ability, Reinforcement} \textendash{} Fokus auf den einzelnen Menschen im Change.
\end{itemize}

\medskip
\noindent\textbf{Arbeitsauftrag:}

\begin{enumerate}
  \item Erstellen Sie eine \textbf{Tabelle} mit den drei Modellen und je drei \emph{St\"arken} und je einem \emph{Schwachpunkt} (z.\,B.\ Kotter: stark in Top-down-Strukturen, schwach bei agilen / lernenden Organisationen).
  \item Ordnen Sie die folgenden Aussagen \emph{einer} der Modell-Phasen / -Schritte zu: ``Das Team versteht, warum wir das tun'' \textperiodcentered{} ``Wir feiern den ersten erfolgreichen Pilot'' \textperiodcentered{} ``Mitarbeiter brauchen ein Training, sonst k\"onnen sie nicht'' \textperiodcentered{} ``Die alte Routine kehrt zur\"uck, sobald der Druck wegf\"allt''.
  \item Welches Modell w\"urden Sie f\"ur eine \emph{technische} Prozess\-ver\"anderung (z.\,B.\ Einf\"uhrung Workflow-Tool) nutzen, welches f\"ur eine \emph{kulturelle} Ver\"anderung (z.\,B.\ ``wir reden offen \"uber Fehler'')? Begr\"unden Sie in je einem Satz.
\end{enumerate}

\ytquery{change management models Kotter Lewin ADKAR comparison strengths weaknesses}

\antwortlinien{14}

\clearpage

% --- AUFGABE 3 -----------------------------------------------------------
\aufgabenkopf{3}{F\"unf Widerstandstypen lesen}{\nivLi}{8\,min}

\noindent Widerstand ist nicht \emph{ein} Ph\"anomen, sondern \emph{mehrere}. F\"unf typische Auspr\"agungen:

\begin{itemize}
  \item \emph{Interessen-Widerstand:} ``Ich verliere Macht / Budget / Sichtbarkeit.''
  \item \emph{Risiko-Widerstand:} ``Ich sehe Gefahren, die niemand adressiert.''
  \item \emph{Identit\"ats-Widerstand:} ``Das war immer meine Rolle / mein Selbstverst\"andnis.''
  \item \emph{\"Uberlastungs-Widerstand:} ``Schon wieder ein Projekt, ich habe keine Zeit.''
  \item \emph{Vertrauens-Widerstand:} ``Ich glaube nicht, dass die F\"uhrung das ernst meint.''
\end{itemize}

\medskip
\noindent\textbf{Arbeitsauftrag:}

\begin{enumerate}
  \item Schreiben Sie f\"ur jeden Typ ein \emph{typisches Zitat aus dem Alltag} (anders als oben).
  \item Notieren Sie f\"ur jeden Typ \emph{eine angemessene Erstreaktion} (z.\,B.\ Interessen: explizit ansprechen + Ausgleichsangebot; \"Uberlastung: Priorit\"aten reduzieren + Sponsor-Backing).
  \item Nennen Sie \textbf{eine Reaktion, die in allen f\"unf F\"allen kontraproduktiv ist} (z.\,B.\ ``Das ist halt jetzt so, bitte ziehen Sie mit'').
\end{enumerate}

\ytquery{change management resistance types interests identity overload trust}

\antwortlinien{12}

\clearpage

% =========================================================================
% L2 AUFGABEN
% =========================================================================
\section*{Niveau L2 \textendash{} Anwendung \& Transfer}
\addcontentsline{toc}{section}{L2 -- Anwendung}

% --- AUFGABE 4 -----------------------------------------------------------
\aufgabenkopf{4}{Change-Architektur f\"ur E2E-Prozessdigitalisierung}{\nivLii}{30\,min}

\begin{infobox}[Ausgangslage]
Ein E2E-Prozess (Order-to-Cash) wird digitalisiert: Workflow-Tool, KPI-Steuerung mit Dashboard, Abweichungsmanagement. Fachbereiche f\"uhlen sich \emph{\"uberwacht} (``KPI = Kontrolle''), IT ist \"uberlastet, Compliance will mehr Kontrollen, Operations will Speed.

\smallskip
\textbf{Ziel:} $-$\,25\,\% Durchlaufzeit, $-$\,30\,\% Rework in 12 Wochen.
\end{infobox}

\noindent\textbf{Arbeitsauftrag:}

\begin{enumerate}
  \item \textbf{F\"unf Rollensteckbriefe} (Sponsor, Change Lead, Process Owner, Line Manager, Champions): f\"ur jede Rolle \emph{drei} konkrete Verantwortungen \emph{f\"ur dieses Vorhaben} (nicht generisch).
  \item \textbf{Risikoliste} (mindestens acht Risiken): z.\,B.\ Akzeptanzverlust, \"Uberlastung, Datenl\"ucken, KPI-Gaming, Schattenprozesse, Compliance-Spannung, IT-Engpass, fehlendes Sponsor-Mandat.
  \item \textbf{Priorisierung}: Top-3 nach \emph{Impact $\times$ Wahrscheinlichkeit} (Matrix 1\,--\,5; kurze Begr\"undung).
  \item \textbf{Zwei Experimente} f\"ur die ersten 2 Wochen: F\"ur welche \emph{zwei Risiken} testen Sie eine Hypothese? Jeweils: \emph{Hypothese}, \emph{Test} (max.\ 2 Wochen), \emph{Erfolg / Misserfolg-Kriterium}.
  \item \textbf{Rituale}: Welche \emph{drei Rituale} verankern den Change in den n\"achsten 12 Wochen (z.\,B.\ Weekly Process Standup 30\,min, Monthly Improvement Board 60\,min, Quarterly Audit-Review)?
\end{enumerate}

\ytquery{change architecture process digitalization roles risks experiments rituals}

\antwortlinien{24}

\clearpage

% --- AUFGABE 5 -----------------------------------------------------------
\aufgabenkopf{5}{Stakeholder-Matrix 2$\times$2 + Einbindungsstrategie}{\nivLii}{25\,min}

\begin{infobox}[Acht Stakeholder]
F\"ur das Vorhaben aus Aufgabe 4 m\"ussen folgende acht Stakeholder eingebunden werden:

\smallskip
\emph{Operations}, \emph{Fachbereich (Vertrieb)}, \emph{IT}, \emph{Compliance}, \emph{Controlling}, \emph{Customer Service}, \emph{Betriebsrat} (falls relevant), \emph{Key Customer (extern)}.
\end{infobox}

\noindent\textbf{Arbeitsauftrag:}

\begin{enumerate}
  \item Erstellen Sie eine \textbf{2$\times$2-Matrix} \emph{Einfluss $\times$ Betroffenheit} (jeweils niedrig / hoch). Ordnen Sie alle acht Stakeholder ein.
  \item Definieren Sie f\"ur jede der vier Felder die \emph{Einbindungsstrategie} aus dem Set: \emph{Informieren} \textperiodcentered{} \emph{Konsultieren} \textperiodcentered{} \emph{Mitentscheiden} \textperiodcentered{} \emph{Verantworten}.
  \item Schreiben Sie f\"ur \emph{drei} ausgew\"ahlte Stakeholder einen \textbf{Steckbrief} (3\,--\,4 Zeilen): \emph{Interessen / Sorgen}, \emph{erwartete Botschaft}, \emph{Kanal}, \emph{Frequenz}.
  \item Identifizieren Sie \textbf{eine ``politische Falle''}: welcher Stakeholder ist scheinbar wenig betroffen, hat aber hohen Einfluss \textendash{} und welche Folge h\"atte es, diesen Stakeholder zu \"ubergehen?
\end{enumerate}

\ytquery{stakeholder mapping matrix influence interest engagement strategy}

\antwortlinien{20}

\clearpage

% --- AUFGABE 6 -----------------------------------------------------------
\aufgabenkopf{6}{Kommunikationspaket: Kick-off + Kernbotschaften}{\nivLii}{25\,min}

\begin{infobox}[Drei blockierende Stakeholder-Gruppen]
\textbf{Team A (Operations):} ``Schon wieder ein Projekt, wir haben keine Zeit.''\\
\textbf{Team B (Fachbereich):} ``KPI-Steuerung\,=\,Kontrolle, wir verlieren Autonomie.''\\
\textbf{Team C (IT):} ``Ohne klare Anforderungen und Datenbasis wird das nichts.''
\end{infobox}

\noindent\textbf{Arbeitsauftrag:}

\begin{enumerate}
  \item \textbf{Kick-off-Botschaft (max.\ 5 S\"atze):} Beginnt mit dem ``Warum'', adressiert die drei Sorgen \emph{namentlich}, kommt zum konkreten Nutzen, benennt eine bewusste Konsequenz (``KPIs sind zur Verbesserung, nicht zur Bestrafung''), und schlie{\ss}t mit einer Einladung.
  \item \textbf{Drei Kernbotschaften} f\"ur die n\"achsten 12 Wochen (jede Botschaft = max.\ 1 Satz; jede adressiert \emph{eine} der drei Sorgen).
  \item \textbf{Zwei Kan\"ale} mit Frequenz und Owner (z.\,B.\ w\"ochentliche Stand-up-Mail durch Change Lead; monatliches Town Hall durch Sponsor; Slack/Teams-Channel mit moderiertem Q\&A).
  \item \textbf{Feedbackmechanismus:} Wie holen Sie regelm\"a{\ss}ig ungefilterte R\"uckmeldung (z.\,B.\ anonyme Pulse Survey 5 Fragen, monatlich; offene Sprechstunde mit dem Sponsor; Champion-Roundtable; Skip-Level-Gespr\"ach)?
  \item \textbf{Was bewusst \emph{nicht} kommunizieren?} Welche zwei Botschaften w\"aren in dieser Phase kontraproduktiv (z.\,B.\ ``Es wird keine Stellenstreichungen geben'' \textendash{} wenn das nicht gesichert ist, zerst\"ort die Aussage Vertrauen)?
\end{enumerate}

\ytquery{change communication kick off key messages channels feedback pulse survey}

\antwortlinien{22}

\clearpage

% --- AUFGABE 7 -----------------------------------------------------------
\aufgabenkopf{7}{Sustainment-Design: KPI-Set + Definition of Done}{\nivLii}{25\,min}

\begin{infobox}[Aufgabenstellung]
Damit der E2E-Prozess nach dem Go-Live \emph{nicht} in die alten Routinen zur\"uckkippt, brauchen Sie ein \emph{Sustainment-Design}: KPIs, Review-Mechanik, Definition of Done, Trainings- und Coachingplan.
\end{infobox}

\noindent\textbf{Arbeitsauftrag:}

\begin{enumerate}
  \item \textbf{KPI-Set (mindestens f\"unf KPIs) \"uber drei Ebenen:}
    \begin{itemize}
      \item \emph{Outcome:} z.\,B.\ Durchlaufzeit, SLA-Quote, Kundenzufriedenheit.
      \item \emph{Prozess:} z.\,B.\ Rework-Quote, Exception-Rate.
      \item \emph{Adoption:} z.\,B.\ \% F\"alle, die im Workflow laufen; Trainingsabschlussquote.
    \end{itemize}
    Pro KPI: \emph{Ziel}, \emph{Owner}, \emph{Datenquelle}, \emph{Review-Frequenz}.
  \item \textbf{Review Board:} Wer sitzt drin, Taktung, Entscheidungsschwellen, Protokoll-Pflichtfelder.
  \item \textbf{Definition of Done (DoD) f\"ur Prozess\"anderungen} \textendash{} f\"unf Kriterien (z.\,B.\ Modell aktualisiert; Owner benannt; KPI angepasst; Training durchgef\"uhrt; Monitoring aktiv).
  \item \textbf{Trainings- / Coaching-Plan} in drei Stufen (Operator \textrightarrow{} Owner \textrightarrow{} Champion) mit Aufwand und Erfolgsnachweis.
  \item \textbf{Anti-Gaming-Mechanik:} Wie verhindern Sie, dass KPIs ``ged\"uckt'' werden (z.\,B.\ Antwortzeiten manipuliert, F\"alle umetikettiert)?
\end{enumerate}

\ytquery{sustainment KPI outcome process adoption review board definition of done}

\antwortlinien{22}

\clearpage

% --- AUFGABE 8 -----------------------------------------------------------
\aufgabenkopf{8}{Fallstudie ``ServiceCo'' \textendash{} E2E-Prozess nachhaltig verankern}{\nivLii}{40\,min}

\begin{infobox}[Fallstudie: ServiceCo]
\textbf{Unternehmen:} ``ServiceCo'' \textendash{} mehrere Standorte; ein E2E-Prozess (Order-to-Delivery) wird workflow- und KPI-gest\"utzt eingef\"uhrt; Ziel $-$\,25\,\% Durchlaufzeit, $-$\,30\,\% Rework in 12 Wochen.

\smallskip
\textbf{Restriktionen:} Zeitdruck; parallele Projekte; unvollst\"andige Daten; Widerstand gegen KPI-Transparenz; Budget begrenzt.
\end{infobox}

\noindent\textbf{Arbeitsauftrag:} Liefern Sie ein integriertes \emph{Change-Paket}, das alle vorhergehenden Bausteine zusammenf\"uhrt:

\begin{enumerate}
  \item \textbf{Change-Architektur:} Rollen (Sponsor / Owner / Lead / Line / Champions); Entscheidungswege; Taktung (Rituale).
  \item \textbf{Stakeholderstrategie:} Acht Stakeholder; 2$\times$2 Matrix; Einbindung pro Quadrant.
  \item \textbf{Kommunikationspaket:} Kick-off-Botschaft, drei Kernbotschaften, zwei Kan\"ale, ein Feedback\-mechanismus.
  \item \textbf{Sustainment-Design:} KPI-Set (Outcome / Prozess / Adoption), Review Board, Trainings- / Coachingplan, Definition of Done.
  \item \textbf{Entscheidung unter Unsicherheit (zwei Trade-off-Protokolle):}
    \begin{itemize}
      \item \emph{Speed vs.\ Compliance:} MVP mit Ausnahmeprozess + Audit Trail? Welche Guardrails sind nicht verhandelbar?
      \item \emph{Transparenz vs.\ Angst:} Wie viel KPI-Transparenz auf Team- vs.\ Personenebene? Welche Guardrails sch\"utzen Kultur?
    \end{itemize}
    Pro Protokoll: \emph{Entscheidung}, \emph{Annahme}, \emph{verworfene Alternative}, \emph{Fr\"uhwarnsignal}, \emph{Reviewdatum}.
\end{enumerate}

\ytquery{change leadership E2E process digital transformation stakeholder sustainment KPI}

\antwortlinien{32}

\clearpage

% =========================================================================
% L3 AUFGABEN
% =========================================================================
\section*{Niveau L3 \textendash{} Management \& Entscheidungsarchitektur}
\addcontentsline{toc}{section}{L3 -- Management}

% --- AUFGABE 9 -----------------------------------------------------------
\aufgabenkopf{9}{Change-Portfolio: WIP-Limits f\"ur Ver\"anderung}{\nivLiii}{30\,min}

\begin{infobox}[Rolle: Chief Transformation Officer]
\textbf{Ihre Rolle:} \emph{CTO (Chief Transformation Officer)} eines Konzerns mit aktuell 17 laufenden Change-Initiativen \"uber 6 Bereiche; Kapazit\"at f\"ur sauber gef\"uhrte Changes liegt realistisch bei 6\,--\,8 parallel.
\end{infobox}

\noindent\textbf{Arbeitsauftrag:}

\begin{enumerate}
  \item \textbf{Priorisierungskriterien (f\"unf):} z.\,B.\ \emph{Strategiebeitrag}, \emph{Dringlichkeit}, \emph{Aufwand}, \emph{Risiko bei Verz\"ogerung}, \emph{Reife der Vorbereitung}. Pro Kriterium ein konkretes Ma{\ss}.
  \item \textbf{WIP-Limit definieren:} Wie viele parallele Changes erlauben Sie pro Bereich / pro Standort / pro Stakeholder-Gruppe? Begr\"undung.
  \item \textbf{Stopp- / Schwenk-Kriterien:} Drei Signale, die ein Change-Vorhaben in den \emph{Pause}- oder \emph{Stop}-Modus heben (z.\,B.\ Adoption nach 8 Wochen $<$\,30\,\%; Sponsor wechselt; KPI-Gaming nachgewiesen).
  \item \textbf{Eskalations-Architektur:} Wer entscheidet \"uber Stopp (Change Lead allein? Steering? Vorstand?), und welche Eskalations\-fristen gelten?
  \item \textbf{Anti-Muster:} Drei typische Fehlmuster in Change-Portfolios (z.\,B.\ ``alles parallel''; ``Sponsor vergibt Mandat ohne Ressourcen''; ``Champions \"uberlasten ohne Anerkennung'') \textendash{} Gegenma{\ss}nahmen je Muster.
\end{enumerate}

\ytquery{change portfolio WIP limits prioritization stop criteria escalation transformation}

\antwortlinien{22}

\clearpage

% --- AUFGABE 10 ----------------------------------------------------------
\aufgabenkopf{10}{Transparenz vs.\ Angst \textendash{} KPI-Guardrails}{\nivLiii}{30\,min}

\noindent\textbf{Diskussionsaufgabe.}

\noindent KPI-Transparenz ist ein zweischneidiges Schwert: \emph{Sichtbar} sein erh\"oht Lerngeschwindigkeit und Verantwortung, \emph{sichtbar gemacht} werden kann aber Angst und Defensivverhalten ausl\"osen. Die Aufgabe der F\"uhrung ist es, eine Linie zu ziehen \textendash{} nicht ``alles transparent'' und nicht ``alles geheim''.

\medskip
\noindent\textbf{Arbeitsauftrag:}

\begin{enumerate}
  \item \textbf{Vier Ebenen} \emph{KPI-Aufl\"osung}: \emph{Unternehmens-aggregiert} \textperiodcentered{} \emph{Bereich} \textperiodcentered{} \emph{Team} \textperiodcentered{} \emph{Person}. F\"ur jede Ebene: \emph{Wer sieht was}, \emph{Zweck}, \emph{Risiko bei Verletzung der Grenze}.
  \item \textbf{Drei Guardrails f\"ur Transparenz}: z.\,B.\ ``Personen-KPIs nur bei begr\"undetem Incident-Fall \textendash{} und mit Vier-Augen''; ``Team-KPIs sind f\"ur Lernen, nicht f\"ur Vergleich''; ``Aggregierte KPIs werden im Townhall geteilt; Bereichs-KPIs in der Bereichsrunde''.
  \item \textbf{Drei KPI-Gaming-Muster}: z.\,B.\ ``Ticket-Schlie{\ss}ungen vor Resolution markieren''; ``Schwierige F\"alle umetikettieren''; ``Antwortzeiten manuell stoppen''. Pro Muster: \emph{Erkennungssignal} und \emph{Gegenma{\ss}nahme}.
  \item \textbf{Kultur-Frage:} Welche zwei Routinen verhindern, dass KPIs als \emph{Druckmittel} eingesetzt werden (z.\,B.\ Retro-Mechanik, ``Was hat geholfen, was nicht?'' statt ``Wer hat versagt?'')?
  \item \textbf{Eine Fr\"uhwarnsignal-Logik}, dass ein KPI in eine Angstkultur umkippt (z.\,B.\ Krankenstand steigt, Tickets ``verschwinden'', offene R\"uckmeldungen werden seltener).
\end{enumerate}

\ytquery{KPI transparency culture fear gaming guardrails psychological safety change}

\antwortlinien{22}

\clearpage

% --- AUFGABE 11 ----------------------------------------------------------
\aufgabenkopf{11}{Transferprojekt \textendash{} Change Leadership Operating Model}{\nivLiii}{50\,min}

\begin{infobox}[Rolle und Ausgangslage]
\textbf{Ihre Rolle:} \emph{Chief Transformation Officer} / \emph{Bereichsleitung} / \emph{Lead Consultant} \textendash{} verantwortlich f\"ur das organisationale Change-System (mehrere parallele Initiativen).

\smallskip
\textbf{Auftrag:} Ein nachhaltiges \emph{Change-Operating-Model} f\"ur Prozessver\"anderungen etablieren, das Geschwindigkeit liefert, Widerstand produktiv macht und Wirkung messbar verankert \textendash{} \"uber mehrere Initiativen.

\smallskip
\textbf{Restriktionen:} Mehrere parallele Changes, begrenzte Kapazit\"at, Stakeholderkonflikte, Datenlage unvollst\"andig, Zeitdruck durch Business-Ziele.
\end{infobox}

\noindent\textbf{Arbeitsauftrag:} Entwickeln Sie eine \emph{Entscheidungsarchitektur} in 8 Schritten:

\begin{enumerate}
  \item \textbf{Top-10 Change-Entscheidungen}: Priorisierung, Ressourcen, KPI-Transparenz-Regeln, Ausnahmeprozess, Eskalationspfade, Kommunikationsregeln, Trainingsinvest, Stoppkriterien, WIP-Limits, Sustainment.
  \item \textbf{Portfolio-Steuerung:} Priorisierungslogik, Taktung, Entlastungsmechanik (WIP-Limits f\"ur Change \textendash{} siehe Aufgabe 9).
  \item \textbf{Stakeholder- \& Kommunikationssystem:} Zielgruppen, Cadence, Feedbackschleifen (Pulse), \emph{Truth Signals} (wie erkennt man \emph{echte} Adoption: Nutzung, Eigeninitiative, Verbesserungs\-vorschl\"age)?
  \item \textbf{Sustainment-System:} DoD-Standards, Review Boards, KPI-Governance, Coaching- / Enablement-Programm, Champions-Netzwerk.
  \item \textbf{Risikomodell:} Widerstandstypen, \"Uberlastungsindikatoren, KPI-Gaming-Erkennung, Gegen\-ma{\ss}nahmen.
  \item \textbf{Roadmap (12 Wochen):} Pilot \textrightarrow{} Skalierung \textrightarrow{} Verstetigung; Messkonzept.
  \item \textbf{Zwei Entscheidungen unter Unsicherheit:}
    \begin{itemize}
      \item \emph{Transparenzregel:} Welche KPI-Transparenz ist hilfreich, welche schadet Kultur? Guardrails + Fr\"uhwarn\-signale.
      \item \emph{Stopp- / Schwenk-Kriterien:} Wann wird ein Change-Vorhaben gestoppt oder neu geschnitten? (WIP, Wirkung, Risiko.)
    \end{itemize}
  \item \textbf{Anschluss an Strategie:} Drei Argumente, warum dieses Operating Model Strategie tr\"agt (Time-to-Effect, Mitarbeiterbindung, Compliance).
\end{enumerate}

\noindent\textbf{Bewertungskriterien (Senior-Level):} Pragmatismus unter Restriktionen \textperiodcentered{} Reife der Trade-off-Logik \textperiodcentered{} kulturelle Sensibilit\"at \textperiodcentered{} Wirkung: \emph{Change wird gef\"uhrt, nicht gemanagt}.

\ytquery{change operating model transformation portfolio sustainment KPI culture}

\antwortlinien{38}

\clearpage

\section*{Abschluss}
\thispagestyle{plain}

\noindent Sie haben alle elf Aufgaben zu Tag 15 bearbeitet. Vergleichen Sie nun Ihre Antworten mit dem \emph{L\"osungsheft Tag 15}.

\medskip
\begin{infobox}[Was Sie jetzt k\"onnen sollten]
\begin{itemize}
  \item Change-Rollen klar besetzen \textendash{} insbesondere die Sponsor-Rolle.
  \item Lewin, Kotter und ADKAR als komplement\"are Werkzeuge einsetzen.
  \item Widerstand als \emph{Information} lesen \textendash{} nach Typ unterschiedlich reagieren.
  \item Stakeholder \"uber Einfluss / Betroffenheit einbinden und ``politische Fallen'' erkennen.
  \item Kommunikationspakete mit glaubw\"urdiger Kernbotschaft entwickeln und Feedback-Mechaniken einbauen.
  \item Sustainment-Design (KPIs, Review, DoD, Training) bauen, das Refreeze tats\"achlich tr\"agt.
  \item Change-Portfolio mit WIP-Limits und Stoppkriterien f\"uhren.
  \item KPI-Transparenz mit Guardrails so dimensionieren, dass Lernen statt Angst entsteht.
  \item Ein Change-Operating-Model als Entscheidungsarchitektur formulieren.
\end{itemize}
\end{infobox}

\vspace{1cm}
\noindent{\color{lpAccent}\rule{\linewidth}{0.6pt}}\\[4pt]
{\footnotesize\sffamily\color{lpGray}Lernplatt \textperiodcentered{} by Can Siebert \textperiodcentered{} Kurs 71 (Dig 42) \textperiodcentered{} Tag 15 \textperiodcentered{} Aufgabenheft \textperiodcentered{} \textcopyright{} 2026}

\end{document}
